the second.

\clearpage
\setlength{\parskip}{0.55em}
\part{Resource-accounted analytic underpinnings}
\label{part:accounted}

\section{The paper's machine model and its intensional obligations}
\label{sec:paper-machines}

\subsection{Machines, inputs, time, and descriptions}

We first repeat the conventions on which the later change of description
language depends.  This is necessary because the change is intensional: it
must preserve programs, bounds, and programs which inspect other programs,
not only partial functions.

\begin{reminder}
Section~\ref{sec:refresher-tm} distinguished the finite transition table from
its run.  In this part, \(\overline M\) is that table's canonical bit string,
\(|M|\) is its bit length, and \(\Time_{M,x}\) counts transitions of the run.
\end{reminder}

\begin{definition}[The paper's Turing-machine convention]
\label{def:paper-tm}
A \(k\)-input Turing machine has \(k\) input tapes, one work tape,
and one output tape.  Tapes are one-way infinite, indexed by \(\N\), and use
\(\{0,1,\sqcup\}\).  Work and output tapes are initially blank.  If the
machine halts, its output is the longest nonblank initial string on the output
tape, with the all-blank output interpreted as the one-bit string \(0\).  On
\(x=(x_1,\ldots,x_k)\), write \(M(x)\) for this string and write \(M(x)=\bot\)
if it never halts.  Its instance time
\(\Time_{M,x}\in\N\cup\{\infty\}\) is the number of transitions to the halt
state.

Every machine has canonical binary code \(\overline M\).  The code lists its
finite states and transition table using a fixed reasonable encoding; its bit
length is \(|\overline M|\), abbreviated \(|M|\).  For a bit string
\(\alpha\), \([\alpha]_k\) is the \(k\)-input machine described by \(\alpha\)
(with the fixed convention for malformed strings).  Tuples are encoded by
\(\operatorname{enc}_k\): each data bit is replaced by
\(0\mapsto01\), \(1\mapsto10\), and each component ends in \(00\).  Thus
\(|\operatorname{enc}_k(x)|=2\sum_i|x_i|+2k\), and encoding takes linear time.
\end{definition}

This is the model in the Turing-machine subsection of the paper.
\gt{sec:tms}{ground-truth/gt-03-prelim.tex}  Its universal-machine convention
is quantitative.  For each \(k\) there is a fixed one-tape \(U_k\), and
universal constants \(C_U,c_U\geq1\), such that if \([\alpha]_k(x)\) halts in
\(T\) steps then \(U_k(\operatorname{enc}_{k+1}(\alpha,x))\) produces the same
answer within
\[
 C_U\bigl(k|\alpha||x|T\bigr)^{c_U},
 \qquad |x|:=\sum_i|x_i|.
\]
The paper also assumes that a high-level machine which invokes a fixed
submachine has description length linear in the embedded submachine code.
Hardwiring a string \(z\) into one input of \(M\) is an effective source
transformation \(s(M,z)\); its output description has length
\(O(|M|+|z|+1)\), is produced in time \(O(|M|+|z|+1)\), and incurs polynomial
simulation overhead.  A timeout wrapper counts down in binary: before the
one-tape simulation it uses \(O(T\log T)\) time and \(O(|M|)\) description
length.  These are conventions, not exact equalities between particular
machine encodings.

\begin{definition}[Decider]
\label{def:paper-decider}
A decider is a five-input machine \(D\) which halts with one bit on every
\((n,x,y,a,b)\), where \(n\) is a positive integer in binary and
\(x,y,a,b\in\{0,1\}^*\).  The first input is the \emph{index}, and
\[
 \Time_D(n)=\max_{x,y,a,b}\Time_{D,(n,x,y,a,b)}.
\]
The maximum is understood in \(\N\cup\{\infty\}\); a finite bound means that
the machine cannot spend time reading arbitrarily remote symbols of the four
unrestricted input tapes.  Output \(1\) is acceptance and \(0\) is rejection.
\end{definition}

This is exactly \code{def:decider}.
\gt{def:decider}{ground-truth/gt-05-games-normalform.tex}

\begin{definition}[Conditionally linear sampler]
\label{def:paper-sampler}
A sampler \(S\) is a six-input Turing machine.  On index \(n\) its legal modes
are
\[
\begin{array}{ll}
(n,\mathsf{dimension}),&
(n,w,\mathsf{marginal},j,z),\\
(n,w,\mathsf{linear},j,u,y),&
(n,w,\mathsf{factor},j,u),
\end{array}
\]
where \(w\in\{\mathrm A,\mathrm B\}\).  Its outputs must describe the ambient
dimension, the \(j\)-th marginal, the conditional linear map, and the
coordinate indicator of its factor space for a pair of level-\(\ell\) CL
functions over \(\F_{q(n)}^{s(n)}\).  The paper writes \(\Time_S(n)\) for the
number of steps before \(S\) halts on index \(n\), i.e. the worst legal call
at that index.  Inputs and outputs which denote integers, finite-field
elements, vectors, and modes are represented by fixed binary conventions.
\end{definition}

This transcribes \code{def:sampler}; its four interfaces and the fact that no
call uses more than six tapes are explicit in the source.
\gt{def:sampler}{ground-truth/gt-04-cl.tex}

\begin{definition}[Normal form and \(\lambda\)-boundedness]
\label{def:paper-lambda}
A normal-form verifier is \(V=(S,D)\), where \(S\) is a CL sampler with field
size \(q(n)=2\) and \(D\) is a decider.  Its description length and level are
\[
 |V|=\max\{|S|,|D|\},\qquad
 \operatorname{level}(V)=\operatorname{level}(S).
\]
For an integer \(\lambda\geq1\), \(V\) is \(\lambda\)-bounded exactly when
\[
 |V|\leq\lambda,qquad
 \Time_S(n)\leq n^\lambda,qquad
 \Time_D(n)\leq n^\lambda\quad(n\geq2).
\]
No bound is required at \(n=1\).
\end{definition}

These are \code{def:normal-ver} and \code{def:lambda}.
\gt{def:normal-ver, def:lambda}{ground-truth/gt-05-games-normalform.tex}
In the game \(V_n\), question and answer alphabets are padded to lengths
\(\Time_S(n)\) and \(\Time_D(n)\), respectively.  The maximum-over-all-inputs
definition above is therefore stronger than a promise only about strings of
those lengths.

\subsection{What a verifier description means}

\begin{reminder}
A description is syntax as data, not an oracle for behavior.  The distinction
was illustrated by the code of the equality machine in
Section~\ref{sec:refresher-tm}; it is now used in anger because compilers scan,
copy, and hardwire verifier code.
\end{reminder}

A ``description of a verifier'' is the ordered pair
\((\overline S,\overline D)\) of canonical machine strings.  It is neither an
oracle for the functions computed by \(S,D\) nor an equivalence class of
programs.  An algorithm receiving it may count its bits, copy it, reject it as
malformed, put it on the input tape of a universal simulator, and manufacture
a new string with parts of it hardwired.  Two extensionally equal deciders can
have different \(|D|\), running times, and behavior under those source
transformations.

The compression theorem makes this interface literal.  There is a
polynomial-time machine \(\Compress\) whose input is
\((\overline V,\lambda)\), with \(\overline V=(\overline S,\overline D)\), and
whose output is the description of a nine-level verifier.  It calls, in order,
\begin{align*}
 \overline{V^{(1)}}&=\ComputeIntroVerifier(\overline V,\lambda,9),\\
 \overline{V^{(2)}}&=\ComputeAnsVerifier
      (\overline{V^{(1)}},\lambda,\mu,\gamma),\\
 \overline{V^{(3)}}&=\ComputeParrepVerifier
      (\overline{V^{(2)}},\lambda,\tau).
\end{align*}
Each call returns another pair of machine descriptions.
\gt{fig:compress, thm:compression}{ground-truth/gt-12-compression.tex}
Likewise, \code{thm:ar} states that \(\ComputeAnsVerifier\) ``outputs the
description of the answer reduced verifier,'' and its proof separately builds
the sampler and decider descriptions.
\gt{thm:ar}{ground-truth/gt-10-answer-reduction.tex}

\subsection{Universal and source-level obligations}

\begin{reminder}
Universality and \(s\)-\(m\)-\(n\), reviewed in
Section~\ref{sec:refresher-universality}, justify calling a supplied program
and fixing some of its arguments.  Every use below must additionally pay for
the universal simulation and for printing the specialized description.
\end{reminder}

The following list is the bookkeeping that a homoiconic description language
must discharge.  Calling the syntax ``lambda calculus'' removes none of these
obligations.

\begin{enumerate}
\item \textbf{Simulation and timeouts.}  Every high-level instruction which
calls \(S\), \(D\), or a generated verifier uses a universal machine, with the
polynomial overhead and timeout counter from \code{sec:tms}.  Hardwiring must
produce finite code with the asserted length.

\item \textbf{Introspection.}  The introspection decider contains the
description of the original sampler and invokes that sampler to obtain its
dimension, marginals, linear maps, and factor spaces.  The source constructs a
constant-size raw universal decider and then hardwires
\((\overline V,\lambda,\ell)\); it first scans \(|V|\) and replaces an
overlong description by a trivial verifier.  Thus introspection consumes the
sampler description, not merely its induced distribution.
\gt{lem:intro-decider-complexity, thm:introspection}{ground-truth/gt-08-introspection.tex}

\item \textbf{Answer reduction.}  In Step 1 of \code{fig:pcpverifier}, the
answer-reduced decider computes
\[
 \mathsf{Circuit}=\mathsf{PaddedSuccinctDecider}
   (\overline D,n,T,Q,\sigma,x,y)
\]
from the \emph{description} of \(D\).  The Cook--Levin circuit depends on the
transition table even though it is much smaller than the computation tableau.
The construction uses \(|D|\leq\sigma\) quantitatively.
\gt{fig:pcpverifier, prop:standard-succinct-sat}{ground-truth/gt-10-answer-reduction.tex}

\item \textbf{Compiler composition.}  The three machines displayed in the
preceding subsection traverse and assemble source descriptions.  The proof
that \(\Compress\) is polynomial time uses their running times.  The
independence of the compressed sampler is also a statement about which source
strings are embedded in its output.

\item \textbf{Self-reference.}  In the Halting section, the eight-input
machine \(\mathcal F\) is run as
\[
 \mathcal F(\overline{\mathcal F},\overline M,
             \lambda,n,x,y,a,b).
\]
It constructs the five-input program obtained by hardwiring its first three
arguments, pairs that decider description with a computed sampler description,
and passes the pair to \(\Compress\).  This is a source-level diagonalization,
not recursive invocation alone.
\gt{fig:halt_f}{ground-truth/gt-12-compression.tex}
\end{enumerate}

Sections~\ref{sec:resource-lambda}--\ref{sec:self-reference} construct one
description language which meets precisely these obligations.  The point is
not Church--Turing equivalence; the point is a costed equivalence between
finite, inspectable descriptions.

\section{A lambda calculus with resources}
\label{sec:resource-lambda}

\subsection{Syntax, phases, and canonical bytes}

\begin{reminder}
The bare calculus of Section~\ref{sec:refresher-lambda} had variables,
abstraction, application, beta reduction, and quotation.  The language
\(\mathcal L\) below is a typed, call-by-value engineering IR: de Bruijn
addresses make binding canonical, and explicit fuel makes evaluation cost a
first-class quantity.
\end{reminder}

Fix a finite set of ground sorts including \(\mathsf{Bit}\),
\(\mathsf{Nat}\), \(\mathsf{Bytes}\), \(\mathsf{Tuple}\),
\(\mathsf{MachineDesc}\), and \(\mathsf{Quoted}(A)\).  Function sorts are
finite products followed by a result sort.  The raw terms of the description
language \(\mathcal L\) are
\[
\begin{split}
t ::= {}& \operatorname{BoundVar}(d,j)
 \mid \operatorname{Hole}(h,A)
 \mid \operatorname{Lambda}(r,t)
 \mid \operatorname{Apply}(t,t^*)
 \mid \operatorname{Fix}(t)\\
&\mid \operatorname{If}(t,t,t)
 \mid \operatorname{Prim}(p,B,t^*)
 \mid \operatorname{Quote}(\operatorname{Closed}(t))\\
&\mid \operatorname{Eval}(t_{\rm code},t^*,F)
 \mid \operatorname{Specialize}(t_{\rm code},\rho).
\end{split}
\tag{8.1}\label{eq:l-grammar}
\]
Here \(t^*\) is a finite ordered list.  \(\operatorname{BoundVar}(d,j)\)
addresses slot \(j\) of the environment frame \(d\) binders outward; hence
alpha-conversion is absent from the representation.  A hole is a named,
typed, static substitution site, not a dynamic variable.  We impose the
\emph{affine-hole discipline}: each hole name occurs at most once.  Reuse of
static data is expressed by binding it once and using de Bruijn variables.
\(F\) is either \(\operatorname{FuelLiteral}(m)\) or a closed arithmetic term
\(\operatorname{FuelBound}(t_n,t_\lambda)\).

Each primitive name \(p\) has a total type contract, a deterministic reference
machine, and a natural-valued charge \(\kappa_p(v_1,\ldots,v_r)\).  The
serialized annotation \(B\) proves
\(\kappa_p(v)\leq B(|v_1|,\ldots,|v_r|)\).  The registry contains the usual
finite byte, bit, tuple, list, comparison, and arithmetic operations, as well
as the bounded machine-table operation used in
Theorem~\ref{thm:tm-to-l}.  A primitive is not an oracle: its reference
machine runs in at most
\(c_p(1+\kappa_p+\sum_i|v_i|)^{e_p}\) Turing steps, with \(c_p,e_p\) fixed in
the language definition.

The sorting judgment \(\Gamma;H\vdash t:A\) is the ordinary simply typed
call-by-value judgment for variables, abstraction, application, and
conditionals, extended as follows.  \(H\) records affine holes;
\(\Quote\) requires a scoped closed term with empty \(H\);
\(\Specialize\) requires an environment covering exactly \(H\) with terms of
the declared sorts; \(\Eval\) requires quoted code, arguments matching the
quoted function sort, and natural fuel.  \(\operatorname{Fix}(P)\) requires a partial body
\(P:A\) with the single hole
\(\mathsf{self\_code}:\mathsf{Quoted}(A)\), and closes that hole.  We write
\(P\{A\}\) for sorted syntax, \(\operatorname{Closed}(P)\) for empty variable
and hole contexts, and \(\operatorname{Quoted}\{A\}\) for its canonical bytes.
Raw ill-sorted syntax is retained so that the evaluator can return
\(\mathsf{TypeError}\) rather than throw a host exception.

\begin{definition}[Canonical serialization]
\label{def:l-serialization}
Let \(\nu(m)=1^\ell0b\), where \(b\) is the \(\ell\)-bit binary expansion of
\(m+1\) and \(\ell=\lfloor\log_2(m+1)\rfloor+1\).  Thus
\(|\nu(m)|=2\ell+1\), and \(\nu\) is prefix-free.  Encode a byte string \(z\)
as \(\nu(|z|)z\).  Encode a term by a four-bit constructor tag followed, in
the order displayed in \eqref{eq:l-grammar}, by \(\nu\)-encoded integer and
string fields, the encoded child count, and the recursive child encodings.
Sorts, primitive contracts, fuel expressions, and static environments use the
same rule with fixed tags and lexicographically ordered field names.  No node
contains a redundant subtree length.

The resulting map \(\operatorname{ser}\) is injective and self-delimiting.
For a term or code value \(t\), define \(|t|:=|\operatorname{ser}(t)|\) in
bits.  This is the description length used below.
\end{definition}

Prefix-freeness of \(\nu\) determines the end of every field; arities determine
the number of recursive parses.  Induction on the first constructor tag
therefore gives a unique parse.  In particular, description size is a number
computed from syntax, not the size of a runtime closure.

The operative representations are:
\begin{center}
\begin{tabular}{P{0.19\textwidth}P{0.34\textwidth}P{0.35\textwidth}}
\toprule
Object & Contents & Legitimate consumer \\
\midrule
Closure & Body pointer and runtime environment & Evaluator only; it has no
canonical project byte string \\
Quoted syntax & Canonical bytes of closed \(\mathcal L\) syntax & Evaluator,
specializer, compiler, and \(\Compress\) \\
Circuit & Finite Boolean DAG with named input blocks & Succinct clause
evaluator and Tseitin transformation \\
Universal evaluator & One fixed program interpreting quoted syntax & Timeout
and bounded-trace construction \\
Specialization & Capture-free code-to-code replacement & Source
transformations; its result is syntax, not a closure \\
\bottomrule
\end{tabular}
\end{center}

\subsection{Small-step semantics and fuel}

\begin{reminder}
Normal form alone does not specify the cost of reaching it.  Recall the term
\(K I\Omega\): normal order terminates and call-by-value does not.  We now fix
one deterministic call-by-value transition system, so a fuel count denotes a
particular run rather than an arbitrary beta-reduction path.
\end{reminder}

We use a deterministic call-by-value CEK machine.  A runtime closure is
\([\lambda r.t,\eta]\), where \(\eta=(\eta_0,\eta_1,\ldots)\) is a stack of
frames and \(\eta_d[j]\) supplies \(\operatorname{BoundVar}(d,j)\).  Other
values are canonical ground data and \(\operatorname{Code}(t)\) for closed
syntax.  A configuration is
\[
 C=\langle q,\eta,K,H,f\rangle,
\]
where \(q\) is the current term or value, \(K\) is a continuation stack,
\(H\) is an append-only heap for lists, tuples, closures, and syntax pointers,
and \(f\in\N\) is remaining fuel.  Arguments are evaluated left to right.

For reference, the contractions after all indicated subterms have become
values are the following.  The notation \(C\longrightarrow_k C'\) means that
the transition has charge \(k\); continuations and heaps not shown are carried
unchanged.
\[
\begin{array}{rcll}
\langle\operatorname{BoundVar}(d,j),\eta,K,H,f\rangle
 &\longrightarrow_1&\langle\eta_d[j],\eta,K,H,f-1\rangle
 &\text{if the address exists},\\
\langle\operatorname{Lambda}(r,t),\eta,K,H,f\rangle
 &\longrightarrow_1&\langle[\lambda r.t,\eta],\eta,K,H,f-1\rangle,&\\
\langle\operatorname{Apply}([\lambda r.t,\eta],\vec v),K,H,f\rangle
 &\longrightarrow_1&\langle t,(\vec v)::\eta,K,H,f-1\rangle
 &(|\vec v|=r),\\
\langle\operatorname{If}(1,v_1,v_0),K,H,f\rangle
 &\longrightarrow_1&\langle v_1,K,H,f-1\rangle,&\\
\langle\operatorname{If}(0,v_1,v_0),K,H,f\rangle
 &\longrightarrow_1&\langle v_0,K,H,f-1\rangle,&\\
\langle\operatorname{Prim}(p,B,\vec v),K,H,f\rangle
 &\longrightarrow_{\kappa_p(\vec v)}&
 \langle p(\vec v),K,H,f-\kappa_p(\vec v)\rangle,&\\
\langle\operatorname{Quote}(\operatorname{Closed}(t)),K,H,f\rangle
 &\longrightarrow_1&\langle\operatorname{Code}(t),K,H,f-1\rangle.&
\end{array}
\tag{8.2}\label{eq:cek-core}
\]
An absent variable address, failed primitive contract, nonbit condition,
wrong application arity, or noncode argument to Eval/Specialize has the sole
successor \(\mathsf{TypeError}\) at charge one.  Evaluation contexts supply
the omitted rules: for example, Apply first pushes a frame remembering its
unevaluated arguments, evaluates the operator, then evaluates each argument
from left to right.  Thus the display specifies contractions while the context
discipline specifies their unique order.

All administrative CEK transitions cost one unit.  They include variable
lookup, creating a closure, pushing or popping one continuation frame,
selecting a Boolean branch, and returning a value.  Beta reduction of
\(\operatorname{Apply}([\lambda r.t,\eta],v_1,\ldots,v_r)\) costs one and
continues with \(t\) in environment
\(((v_1,\ldots,v_r),\eta_0,\eta_1,\ldots)\).  A wrong arity, applying a
nonclosure, or using a nonbit condition takes one transition to
\(\mathsf{TypeError}\).  A residual \(\operatorname{Hole}\) does the same.
\(\operatorname{Quote}(\operatorname{Closed}(t))\) takes one transition to the
syntax pointer \(\operatorname{Code}(t)\); serialization was performed when
the enclosing description was admitted, so this transition does not copy
\(|t|\) bits.

After its arguments have become values,
\(\operatorname{Prim}(p,B,v_1,\ldots,v_r)\) costs exactly
\(\kappa_p(v_1,\ldots,v_r)\geq1\).  If the contract or dynamic sorts fail it
uses one unit and returns \(\mathsf{TypeError}\).  For purposes of the local
semantics, a charge \(k\) is expanded into \(k\) deterministic microsteps of
the primitive reference machine; thus every microstep consumes one fuel unit.

\(\operatorname{Specialize}(\operatorname{Code}(P),\rho)\) first checks the
sort and coverage of \(\rho\), then traverses \(P\) in prefix order.  It uses
\(1+|P|+\sum_h|\rho(h)|\) fuel and returns the code value of the substituted
term, or \(\mathsf{TypeError}\).  \(\operatorname{Eval}\) evaluates its code,
arguments, and fuel expression, installs a delimiter holding the requested
inner budget, and runs the same CEK transition relation inside it.  Installing
and removing that delimiter costs two units.  Each inner microstep decrements
both the inner budget and the ambient budget.  Exhausting either produces
\(\mathsf{OutOfFuel}\).

Finally, \(\operatorname{Fix}(P)\) is a value of the result sort of \(P\).
When it is applied to arguments \(u\), one unfolding transition installs the
static binding
\[
 \mathsf{self\_code}\longmapsto
 \operatorname{Code}(\operatorname{Fix}(P))
\]
and evaluates \(P\) with the same arguments.  This is a constant-size heap
link, not a serialization of a recursively expanded tree.  Including the
application and binding frames, an unfolding costs the fixed constant
\(c_Y=3\).  The corresponding fully specialized syntax can be materialized by
\(\Specialize\) when a source consumer asks for it.

If a transition has charge \(k\) and \(f<k\), its unique successor is
\(\mathsf{OutOfFuel}\); otherwise the successor has counter \(f-k\).  A final
configuration is exactly one of
\[
 \mathsf{Value}(v),\qquad \mathsf{OutOfFuel},\qquad
 \mathsf{TypeError}.
\]
We write \(\operatorname{run}(t;f)\) for this result.  For closed function
syntax \(t\) and an argument tuple \(u\),
\(\operatorname{eval}_{\mathcal L}(t,u;f)\) abbreviates the run of
\(\operatorname{Apply}(t,\overline u)\), where \(\overline u\) is the literal
value syntax.  Results do not record unused fuel.

For the paper, the next theorem buys a unique bounded-evaluation trace for
each quoted verifier call, which can later be encoded by Cook--Levin.
\begin{theorem}[Determinism and outcome trichotomy]
\label{thm:l-determinism}
Every nonfinal configuration with positive fuel has at most one successor.
For every closed raw term and finite initial fuel, evaluation terminates in
exactly one of \(\mathsf{Value}(v)\), \(\mathsf{OutOfFuel}\), or
\(\mathsf{TypeError}\).
\end{theorem}

\begin{proof}
The current control item is either a term, a value, or a primitive microstate.
Constructor tags are disjoint.  For a term tag, the evaluation-context rule
selects the leftmost child not yet represented by a value; if there is one,
exactly one continuation frame is pushed.  If all children are values, exactly
one contraction above applies.  Variable addresses select at most one frame
and slot.  Primitive reference machines and their contracts are deterministic
by definition.  The Boolean and arity tests have disjoint success and error
cases.  The same reasoning applies to the unique top Eval delimiter and to the
single self-code binding of Fix.  Hence two distinct successors cannot exist.

Every nonfinal transition consumes at least one unit.  Starting from fuel
\(f\), at most \(f\) transitions or primitive microsteps can occur before a
value or error is returned; if neither has occurred, the next requested step
returns out-of-fuel.  The three final tags are disjoint, and determinism makes
the reached tag unique.
\end{proof}

For the paper, the next theorem buys safe enlargement of a proved timeout:
extra fuel cannot change a completed verifier answer.
\begin{theorem}[Fuel monotonicity]
\label{thm:fuel-monotonicity}
If \(\operatorname{run}(t;f)=R\) with
\(R\neq\mathsf{OutOfFuel}\), then
\(\operatorname{run}(t;f+g)=R\) for every \(g\geq0\).  More precisely, if the
first run uses \(c\leq f\) units, the second uses the same \(c\) transitions
and leaves \(g+(f-c)\) units.
\end{theorem}

\begin{proof}
Induct on the \(c\) transitions of the terminating run.  At step \(i\), the
larger-fuel run has exactly \(g\) more ambient fuel and, inside each Eval
delimiter, exactly the same inner budget as the smaller run.  The smaller run
can pay the uniquely determined charge \(k_i\), so the larger run can pay it
as well.  Determinism gives the same nonfuel successor data, and subtraction
preserves the difference \(g\).  After \(c\) steps the larger run therefore
reaches the same final tag and value.  No operational rule branches on the
numeric fuel except the insufficient-fuel rule, which was not selected in the
smaller run.
\end{proof}

\subsection{Quotation, evaluation, and specialization}

\begin{reminder}
Quotation turns a finite syntax tree into data; specialization is the lambda
counterpart of \(s\)-\(m\)-\(n\).  Neither operation decides what the quoted
program computes.  It only parses or rewrites its canonical tree.
\end{reminder}

External evaluation of bytes must pay for decoding even though an internal
Quote is a syntax pointer.  Define
\[
 h(d,u)=3+|d|+|\operatorname{enc}(u)|.
\]
The quoted evaluator \(\Eval_{\mathcal L}(d,u;F)\) checks and decodes \(d\),
installs the argument tuple, and then runs the CEK machine.  By definition the
front-end scan uses exactly \(h(d,u)\) fuel; malformed code returns
\(\mathsf{TypeError}\) during that scan.

For the paper, the next theorem buys the right to replace execution of a term
by execution of its transmitted description while charging for decoding.
\begin{theorem}[Quote/Eval equation with explicit overhead]
\label{thm:quote-eval}
Let \(t\) be a closed function term, let \(d=\operatorname{ser}(t)\), and let
\(u\) have the required argument sorts.  For every \(f\geq0\),
\[
 \Eval_{\mathcal L}(\Quote(t),u;f+h(d,u))
   =\operatorname{eval}_{\mathcal L}(t,u;f).
\tag{8.3}\label{eq:quote-eval}
\]
Here the left side takes the canonical bytes denoted by \(\Quote(t)\).  If the
right side uses \(c\leq f\) units and returns a value or type error, the left
side uses exactly \(h(d,u)+c\).  If it requests a \((f+1)\)-st unit, both sides
return out-of-fuel after their respective budgets are exhausted.
\end{theorem}

\begin{proof}
The prefix decoder is inverse to serialization by
Definition~\ref{def:l-serialization}; hence after its fixed scan the left
machine holds the same syntax tree \(t\), literal arguments, empty lexical
environment, and application continuation as the right machine.  The left
counter is then exactly \(f\).  From this point the two configurations are
identical.  Theorem~\ref{thm:l-determinism} gives identical successors until a
final outcome.  The front-end accounts for the stated additional
\(h(d,u)\) units and no other rule differs.
\end{proof}

\begin{lemma}[Capture-free specialization and its size]
\label{lem:specialization}
Let \(P\) be well-scoped partial syntax satisfying the affine-hole discipline,
and let \(\sigma\) map every hole of \(P\) to closed, sort-correct syntax.
There is a unique term \(\Specialize(P,\sigma)\), it has no holes, it preserves
the sort of \(P\), and
\[
 |\Specialize(P,\sigma)|
 \leq |P|+\sum_{h\in\operatorname{dom}\sigma}|\sigma(h)|.
\tag{8.4}\label{eq:specialize-size}
\]
The construction takes at most \(1+|P|+\sum_h|\sigma(h)|\) CEK microsteps.
\end{lemma}

\begin{proof}
Traverse \(P\) in prefix order.  Copy every nonhole tag and payload unchanged.
At \(\operatorname{Hole}(h,A)\), check the unique table entry, check its sort,
and emit \(\sigma(h)\).  The inserted term is closed, so its de Bruijn indices
refer to no binder in \(P\); indices in the surrounding syntax are also
unchanged.  Thus no variable can be captured.  Induction on \(P\) proves sort
preservation, because the hole and replacement have the same sort at the only
new case.  Coverage removes every hole, and deterministic traversal gives
uniqueness.

Serialization has no ancestor subtree-length fields.  Replacing the one
occurrence of \(h\) deletes a nonempty encoded Hole node and inserts exactly
\(|\sigma(h)|\) bits.  Since distinct affine holes have disjoint occurrences,
\[
 |\Specialize(P,\sigma)|
 =|P|-\sum_h|\operatorname{Hole}(h,A_h)|+\sum_h|\sigma(h)|,
\]
which implies \eqref{eq:specialize-size}.  One scan step per input or emitted
bit, plus the initial check, gives the time bound.  Without the affine-hole
discipline the correct last term would sum once per hole occurrence; this is
why the discipline is part of the formal system.
\end{proof}

\section{Simulation theorems: the machine--term bridge}
\label{sec:simulation}

\begin{reminder}
Section~\ref{sec:refresher-engineering} gave the two compiler ideas and warned
that Church--Turing equivalence alone says nothing quantitative.  This section
supplies the description-length and running-time inequalities needed by the
compression theorem.
\end{reminder}

\subsection{From the paper's machines to
\texorpdfstring{\(\mathcal L\)}{L}}

We use explicit functional tapes and one bounded table primitive.  A tape is a
zipper \((L,s,R)\): \(L\) is the reversed list strictly left of the head,
\(s\in\{0,1,\sqcup\}\) is the scanned symbol, and \(R\) is the list strictly
right of it with implicit blank tail.  A right move maps
\((L,s,rR')\) to \((sL,r,R')\), taking \(r=\sqcup\) when \(R\) is empty; a
left move is symmetric.  These are fixed list primitives of charge at most
four.  A configuration is a tuple of the finite control state and the
\(k+2\) zippers for input, work, and output tapes.

For a machine description \(m\), the primitive
\(\mathsf{TMDelta}(m,q,s_1,\ldots,s_{k+2})\) scans the encoded transition table
and returns the next state, written symbols, and head motions.  Its declared
charge is
\[
 \kappa_{\mathsf{TMDelta}}\leq 8(|m|+k+2).
\]
The primitive validates \(m\); malformed descriptions use a fixed rejecting
transition.  Its implementation is an ordinary finite table scan.  Thus the
description of the table remains data of length \(|m|\), rather than being
duplicated into a nest of conditionals.

Let \(\mathsf{init}_k\) parse the \(k\)-tuple encoding into input zippers and
blank work/output zippers, and let \(\mathsf{out}\) read the longest nonblank
output prefix.  Define
\[
\begin{split}
 \mathsf{loop}_m={}&\operatorname{Fix}\bigl(
  \lambda C.\ \operatorname{If}(\mathsf{halt}(C),\mathsf{out}(C),\\[-2pt]
 &\hspace{8em}\mathsf{self}(\mathsf{move}(
       C,\mathsf{TMDelta}(m,\mathsf{view}(C)))))\bigr),\\
 \langle\!\langle M\rangle\!\rangle={}&\lambda z.\mathsf{loop}_{\overline M}
                  (\mathsf{init}_k(z)).
\end{split}
\tag{9.1}\label{eq:tm-interpreter-term}
\]
Here \(\mathsf{self}\) is the closure obtained from the Fix self-code link;
equivalently it is \(\Eval\) of that link with the next configuration.  All
other displayed operations are fixed primitives with the stated list
representation.

For the paper, the next theorem buys a lambda description for every original
sampler or decider without changing its answers or hiding a size blow-up.
\begin{theorem}[TM to \(\mathcal L\), with resources]
\label{thm:tm-to-l}
There are universal constants \(a_0,b_0,c_0\geq1\) such that for every
\(k\)-input machine \(M\), \eqref{eq:tm-interpreter-term} is a closed term and
\[
 |\langle\!\langle M\rangle\!\rangle|\leq a_0|M|+b_0.
\]
If \(M(x_1,\ldots,x_k)\) halts within \(T\geq1\) transitions, then
\[
 \Eval_{\mathcal L}\!\left(
   \Quote(\langle\!\langle M\rangle\!\rangle),
   \langle x_1,\ldots,x_k\rangle;
   c_M T(T+|x|)\right)=M(x_1,\ldots,x_k),
\tag{9.2}\label{eq:tm-to-l-bound}
\]
where \(|x|=\sum_i|x_i|\) and one may take
\(c_M=c_0(|M|+k+1)\).  The tuple on the left is the same dual-rail encoding as
in Definition~\ref{def:paper-tm}.  In particular, a decider receives exactly
\[
 \langle\operatorname{bin}(n),x,y,a,b\rangle_{\mathcal L},
 \qquad
 \operatorname{enc}_{\mathcal L}\bigl(\langle\operatorname{bin}(n),x,y,a,b
   \rangle_{\mathcal L}\bigr)
 =\operatorname{enc}_5(\operatorname{bin}(n),x,y,a,b),
\]
and the five components initialize five separate input zippers.
\end{theorem}

\begin{proof}
The term contains one copy of \(\overline M\), the fixed interpreter body, and
self-delimiting constructor overhead.  Each input bit in \(\overline M\)
contributes a fixed number of serialized bits, so constants \(a_0,b_0\)
independent of \(M\) give the size bound.

Initialization scans the dual-rail input once.  For valid input it consumes at
most \(c_0(|x|+k+1)\) units and produces exactly the paper's initial tapes.
Assume inductively that the functional configuration before loop iteration
\(i\) represents the machine configuration after \(i\) transitions.  The
views return the same scanned symbols.  The table primitive selects the unique
transition of \(M\), and each zipper update writes the selected symbol and
moves in the selected direction.  Therefore the next functional configuration
represents the machine configuration after \(i+1\) transitions.  This proves
the simulation invariant for \(0\leq i\leq T\).

One loop iteration, including the fixed-point unfolding, table scan, tuple
operations, and \(k+2\) zipper updates, uses at most
\(c_0(|M|+k+1)\) units.  No visited zipper has more than \(|x|+T\) explicit
cells.  The final output scan therefore uses at most
\(c_0(|x|+T+1)\) units.  Adding initialization, \(T\) iterations, output, and
the quote/decode overhead from Theorem~\ref{thm:quote-eval} gives at most
\[
 c_0(|M|+k+1)\bigl(|x|+1+T\bigr)+
 c_0(|M|+k+1)T.
\]
For \(T\geq1\), both \(|x|+T+1\) and \(T\) are at most
\(2T(T+|x|)\).  Increasing the universal choice of \(c_0\) by a factor four
makes this no greater than the fuel in \eqref{eq:tm-to-l-bound}.  The invariant
and the definition of \(\mathsf{out}\) then give the same output string.  The
explicit quadratic form is deliberately conservative; the zipper simulation
itself is linear in \(T\) apart from the hardwired table scan.
\end{proof}

\subsection{From \texorpdfstring{\(\mathcal L\)}{L} to the paper's machines}

For the paper, the next theorem buys a route back from the IR to the exact
Turing-machine model in which verifier bounds and compression are stated.
\begin{theorem}[A universal evaluator for \(\mathcal L\)]
\label{thm:l-to-tm}
There is a fixed three-input Turing machine \(U_{\mathcal L}\), of description
length \(b_U=O(1)\), and a universal constant \(C_L\geq1\), such that on input
\((d,u,f)\) it returns the same value, out-of-fuel, or type-error result as
\(\Eval_{\mathcal L}(d,u;f)\), and
\[
 \Time_{U_{\mathcal L}}(d,u,f)
 \leq C_L\bigl(|d|+|u|+f+1\bigr)^6.
\tag{9.3}\label{eq:l-to-tm-bound}
\]
Given closed \(t\), hardwiring \(d=\operatorname{ser}(t)\) produces a machine
\((U_{\mathcal L})_t\) with
\[
 |(U_{\mathcal L})_t|\leq a_U|t|+b_U.
\]
On the paper's separate input tapes it may represent each argument by a lazy
head-and-position handle.  Consequently, a run with fuel \(f\) accesses at
most \(f\) symbols of those tapes and has the refined bound
\(C_L(|t|+f+1)^6\), independent of unread suffixes.  Hence a total decider term
with fuel schedule \(f_D(n)\) is a decider in the paper's sense and
\[
 \Time_D(n)\leq C_L(|t|+f_D(n)+1)^6.
\tag{9.4}\label{eq:decider-term-time}
\]
\end{theorem}

\begin{proof}
The machine first checks the prefix serialization and writes node records on a
work-tape heap.  It stores a node address, environment-frame addresses,
continuation records, and binary fuel counters.  During at most \(f\)
microsteps the heap acquires at most a fixed number of records per step, so its
length is at most \(c(|d|+|u|+f+1)\) for a universal \(c\).  A conservative
single-tape implementation finds a record by a full scan, performs binary
address arithmetic, and returns to the simulated head in at most the square of
that length.  Decoding, \(f\) simulated steps, and final serialization are
therefore bounded by a universal constant times the fourth power of the
length.  Converting the fixed multitape implementation to the paper's
one-tape convention and allowing the paper's polynomial universal-simulation
overhead is bounded by the sixth power after increasing \(C_L\).  This proves
\eqref{eq:l-to-tm-bound}.  Primitive microstates use their fixed reference
machines and have charged running time; their polynomial exponents were fixed
when \(\mathcal L\) was fixed and are included in the exponent six (or, for a
different registry, in another fixed exponent).

The simulator description contains only the parser, CEK transition table,
primitive registry, and serializers, so its length is \(b_U\), independent of
\(d,u,f\).  The paper's hardwiring convention embeds \(d\) once and gives the
linear description bound.  For separate input tapes, a handle stores a tape
number and head position.  One microstep can move such a head by at most the
work performed during that charged step, so no more than \(f\) input symbols
can be inspected.  Removing the untouched encoded suffixes from the preceding
storage account yields the refined bound.  Taking a maximum over all
\((x,y,a,b)\) now proves \eqref{eq:decider-term-time}; arbitrarily long unread
suffixes do not affect it.
\end{proof}

\subsection{The resource dictionary}

\begin{reminder}
The qualitative dictionary in Section~\ref{sec:refresher-engineering} matched
machine time with reduction fuel and machine code with serialized syntax.  The
table below is the quantitative version: it records the polynomials rather
than merely asserting that simulations exist.
\end{reminder}

For a term program, let \(s_t=|t|\) and let \(f_t(n)\) be a proved fuel
schedule which suffices for every legal call at index \(n\).  The translations
above give the following dictionary.

\begin{center}
\small
\begin{tabularx}{\textwidth}{P{0.18\textwidth}P{0.28\textwidth}X}
\toprule
Paper quantity & \(\mathcal L\) counterpart & Quantitative relation \\
\midrule
\(\Time_D(n)\) & decider fuel \(f_D(n)\) & Term to TM:
\(\Time_D(n)\leq C_L(|D_{\mathcal L}|+f_D(n)+1)^6\).  TM to term:
fuel \(c_DT(T+|u|)\). \\
\(\Time_S(n)\) & maximum sampler-mode fuel \(f_S(n)\) & Same two bounds,
taking the maximum over the dimension, marginal, linear, and factor calls. \\
\(|D|\) & \(|\operatorname{ser}(D_{\mathcal L})|\) & Each translation is
linear: at most \(a_0|D|+b_0\) or \(a_U|D_{\mathcal L}|+b_U\). \\
\(|S|\) & \(|\operatorname{ser}(S_{\mathcal L})|\) & The same linear bounds;
static CL tables and source descriptions are counted once. \\
\(\lambda\)-bounded & size at most \(\lambda\), fuel at most \(n^\lambda\) &
Preserved in both directions as \(A\lambda\)-bounded for a universal integer
\(A\), as proved below. \\
Threshold \(n\geq C_0\) & same semantic index condition & Translation does not
change \(n\); resource estimates may replace it by
\(n\geq\max\{C_0,2\}\), never hide it in \(\poly\). \\
\bottomrule
\end{tabularx}
\end{center}

For the paper, the next theorem buys invariance of the verifier class under a
change of programming language, after one universal rescaling of \(\lambda\).
\begin{theorem}[Preservation of \(\lambda\)-boundedness]
\label{thm:lambda-preservation}
Fix the encodings and primitive registry above.  There is a universal integer
\(A\geq1\) such that:
\begin{enumerate}
\item translating any \(\lambda\)-bounded paper verifier term-by-term with
Theorem~\ref{thm:tm-to-l} gives an \(A\lambda\)-bounded
\(\mathcal L\)-verifier;
\item compiling any \(\lambda\)-bounded \(\mathcal L\)-verifier with
Theorem~\ref{thm:l-to-tm} gives an \(A\lambda\)-bounded paper verifier.
\end{enumerate}
The translation preserves the verifier's index and its input/output function.
\end{theorem}

\begin{proof}
For the first direction, \(|M|\leq\lambda\) and the size theorem gives
\(|\langle\!\langle M\rangle\!\rangle|\leq(a_0+b_0)\lambda\), since \(\lambda\geq1\).
At index \(n\geq2\), the paper machine runs for \(T\leq n^\lambda\).  It can
inspect or emit at most \(T\) tape symbols.  Across at most six input tapes the
portion relevant to the computation therefore has encoded length at most
\(8n^\lambda\), after increasing the constant to cover \(\log n\) and mode
tags.  Also
\(c_M\leq c_0(\lambda+7)\leq8c_0n^\lambda\).  Theorem~\ref{thm:tm-to-l}
then requires at most
\[
 8c_0n^\lambda\cdot n^\lambda\cdot9n^\lambda
 =72c_0n^{3\lambda}
 \leq n^{(3+\lceil\log_2(72c_0)\rceil)\lambda}.
\]
The last inequality uses \(n\geq2\) and \(\lambda\geq1\).  It applies to
both the sampler and decider.

For the second direction, hardwiring gives description size at most
\((a_U+b_U)\lambda\).  With fuel at most \(n^\lambda\), the lazy-input form of
\eqref{eq:decider-term-time} is at most
\[
 C_L(\lambda+n^\lambda+1)^6
 \leq C_L(3n^\lambda)^6
 \leq n^{(6+\lceil\log_2(C_L3^6)\rceil)\lambda}.
\]
We used \(\lambda\leq2^\lambda\leq n^\lambda\).  The same calculation holds
for every sampler mode.  Taking \(A\) to be the maximum of the four displayed
description and exponent coefficients proves both assertions.  Semantics is
preserved by Theorems~\ref{thm:tm-to-l} and~\ref{thm:l-to-tm}; neither
translation changes the binary index.
\end{proof}

The constants \(a_0,b_0,c_0,a_U,b_U,C_L\), fixed arities, serialization tags,
and the universal simulation exponent are absorbed when the paper writes
\(O(\cdot)\) or \(\poly(\cdot)\): they are universal after the two languages
are fixed.  The particular \(|D|,|S|\), a supplied \(\lambda\), a fuel or time
parameter, and hardwired \(|M|\) are \emph{arguments} of those polynomials and
cannot be hidden as universal constants.  Nor can the semantic threshold
\(C_0\), the change from \(\lambda\) to \(A\lambda\), or constants appearing
inside completeness and soundness inequalities be erased by \(\poly\)
notation.

\section{Self-reference as a construction on descriptions}
\label{sec:self-reference}

\begin{reminder}
Kleene's theorem in Section~\ref{sec:refresher-kleene} produced a behavioral
fixed point by applying a parameter-fixing compiler to its own code.  We now
repeat that construction with explicit description-size bounds and identify
it with the paper's diagonal verifier.
\end{reminder}

\subsection{The recursion theorem with size accounting}

Let \(\varphi_e^{(k)}\) be the partial function of the \(k\)-input paper
machine with description \(e\).  We use equality of partial functions only in
the conclusion; every operation in the construction acts on strings.

\begin{lemma}[The \(s\)-\(m\)-\(n\) operation]
\label{lem:smn}
For every \(r,k\geq1\) there is a total source transformer \(s_{r,k}\) such
that
\[
 \varphi_{s_{r,k}(e,z_1,\ldots,z_r)}^{(k)}(x)
 =\varphi_e^{(r+k)}(z_1,\ldots,z_r,x).
\]
Under the paper's hardwiring convention there are constants
\(a_s,b_s\), depending only on \(r,k\), for which
\[
 |s_{r,k}(e,z)|\leq a_s\bigl(|e|+\textstyle\sum_i|z_i|\bigr)+b_s,
\]
and the output code is computed in time linear in the right-hand side.
\end{lemma}

\begin{proof}
Output the fixed wrapper which writes
\(\operatorname{enc}_{r+k+1}(e,z_1,\ldots,z_r,x)\) on its work tape and invokes
the universal machine \(U_{r+k}\).  The wrapper contains one copy of each
static bit and constant code for tuple encoding and universal simulation.  It
therefore has the displayed linear length and can be printed in linear time.
The universal-machine theorem gives the stated partial-function equation on
halting inputs; on a nonhalting input both sides simulate the same machine
forever.
\end{proof}

For the paper, the next theorem buys self-reference while keeping the fixed
point's source length linear in the source transformer.
\begin{theorem}[Kleene's second recursion theorem, paper model]
\label{thm:kleene-two}
Let \(Q\) be a total computable transformation from descriptions of
\(k\)-input machines to descriptions of \(k\)-input machines.  There is an
effectively constructible description \(e_Q\) such that
\[
 \varphi_{e_Q}^{(k)}=\varphi_{Q(e_Q)}^{(k)}.
\tag{10.1}\label{eq:kleene-fixed}
\]
If \(q\) is a description of the machine computing \(Q\), then, for constants
\(a_K,b_K\) fixed by \(k\) and the encodings,
\[
 |e_Q|\leq a_K|q|+b_K,
\]
and \(e_Q\) is constructed in time polynomial in \(|q|\).
\end{theorem}

\begin{proof}
By Lemma~\ref{lem:smn}, let \(s(z,z)\) denote the description obtained by
hardwiring two copies of \(z\) into the first two inputs of a suitable
\((k+2)\)-input machine.  Define the partial computable function
\[
 G(z,x)=\varphi_{Q(s(z,z))}^{(k)}(x).
\]
A machine for \(G\) first computes the finite string \(s(z,z)\), runs the total
code transformer \(Q\) on it, and universally evaluates the resulting
\(k\)-input code on \(x\).  Let \(p\) be its description and set
\(e_Q=s(p,p)\).  Then for every \(x\), including the case in which the final
computation diverges,
\[
 \varphi_{e_Q}^{(k)}(x)
 =\varphi_p^{(k+1)}(p,x)
 =\varphi_{Q(s(p,p))}^{(k)}(x)
 =\varphi_{Q(e_Q)}^{(k)}(x),
\]
which proves \eqref{eq:kleene-fixed}.

The program \(p\) contains one copy of \(q\) and fixed code for \(s\) and the
universal evaluator, so \(|p|\leq a_1|q|+b_1\).  Applying the linear size bound
of Lemma~\ref{lem:smn} to \(s(p,p)\) gives
\(|e_Q|\leq2a_s a_1|q|+(2a_sb_1+b_s)\).  These are the asserted
\(a_K,b_K\).  Printing the two wrappers and executing the total hardwiring
transformations takes polynomial time under the paper's conventions.
\end{proof}

\subsection{The call-by-value Z combinator and YCode}

\begin{reminder}
Raw \(Y\) unfolds too early under call-by-value; Section~\ref{sec:refresher-lambda}
derived the eta-delayed \(Z\) combinator.  The constructor below realizes that
delay as a typed, constant-size self-code link.
\end{reminder}

For orientation, call-by-value lambda calculus uses
\[
 Z=\lambda f.(\lambda x.f(\lambda u.xx u))
                 (\lambda x.f(\lambda u.xx u)),
\]
rather than the call-by-name \(Y\), because eager evaluation of \(Yf\) unfolds
before supplying an argument.  This untyped term has constant size, but it
does not itself provide a canonical code value of the fixed point.  Our typed
description constructor does.

For a partial body \(P\) with its one distinguished quoted hole, define
\[
 \operatorname{YCode}(P):=\operatorname{Fix}(P),\qquad
 d_P:=\Quote(\operatorname{YCode}(P)).
\]
The source equation is
\[
 \operatorname{unfold}(\operatorname{YCode}(P))
 =\Specialize(P,\{\mathsf{self\_code}\mapsto d_P\}).
\tag{10.2}\label{eq:ycode-source}
\]
It is an equation between finite descriptions.  Its right side is produced on
demand by Lemma~\ref{lem:specialization}; the runtime evaluator instead uses a
constant-size environment link.

For the paper, the next theorem buys an operational self-code equation whose
unfolding cost and serialized size are both explicit.
\begin{theorem}[Description-level fixed-point equation]
\label{thm:ycode}
For every well-sorted \(P\), argument tuple \(u\), and \(f\geq c_Y=3\),
\[
 \operatorname{eval}_{\mathcal L}(\operatorname{YCode}(P),u;f)
 =\operatorname{eval}_{\mathcal L}\!\left(
   \Specialize(P,\{\mathsf{self\_code}\mapsto d_P\}),u;f-c_Y\right).
\tag{10.3}\label{eq:ycode-eval}
\]
Equivalently, if \(f_P\) is the code transformer represented by \(P\), the
right side is conventionally written
\(\operatorname{eval}_{\mathcal L}(f_P(\operatorname{YCode}(f_P)),u;
f-c_Y)\).  Moreover
\[
 |\operatorname{YCode}(P)|=|P|+c_{\rm fix}
\]
for the fixed four-bit Fix tag and its fixed sort annotation
\(c_{\rm fix}=O(1)\).
\end{theorem}

\begin{proof}
The unique Fix application rule performs three administrative actions: push
the application frame, install the self-code heap link, and enter \(P\).  The
resulting control, ordinary environment, and continuation are the same as for
the specialized body; a lookup of the sole hole follows the installed link and
therefore returns \(d_P\).  Subsequent transitions coincide by
Theorem~\ref{thm:l-determinism}.  Exactly three units have been consumed, which
proves \eqref{eq:ycode-eval}, including matching out-of-fuel and type-error
outcomes.  Canonical serialization of Fix writes its fixed tag and annotation
followed by one copy of the serialization of \(P\); it never expands the
self-code link.  This proves the size equality.
\end{proof}

\subsection{The paper's explicit diagonal program}

\begin{reminder}
The expression \(\mathcal F(\mathcal F,M,L,\ldots)\) is the recursion theorem
unrolled: hardwiring the first copy of \(\mathcal F\) manufactures the
five-input decider whose code is then supplied to \(\Compress\).
\end{reminder}

Let \(M:P\{\mathsf{MachineDesc}\}\) and
\(L:P\{\mathsf{Level}\}\) be closed literal terms, let \(S_L\) be the sampler
description returned by \(\ComputeSampler(L)\), and let \(n,x,y,a,b\) be the
five bound inputs.  The partial body corresponding to \code{fig:halt_f} is
\[
\begin{split}
 \Psi_{M,L}=\lambda(n,x,y,a,b).\ \operatorname{If}(&
  \operatorname{Prim}(\mathsf{halts\_within},M,n),\ \mathsf{true},\\
 &\operatorname{Eval}(
   \operatorname{ans}(\operatorname{Eval}(
      \Quote(\Compress),\\[-2pt]
 &\quad (\operatorname{quoted\_pair}(\Quote(S_L),
       \operatorname{Hole}(\mathsf{self\_code},
          \mathsf{Quoted}(\mathsf{Decider}))),L),F_C)),\\[-2pt]
 &\quad(n,x,y,a,b),\operatorname{FuelBound}(n,L))).
\end{split}
\tag{10.4}\label{eq:psi-ml}
\]
The finite-fuel symbol \(F_C\) is any proved bound for the terminating
\(\Compress\) call; \(\operatorname{ans}\) selects the returned decider
description.  Define
\[
 D_{M,L}=\operatorname{YCode}(\Psi_{M,L}),\qquad
 V_{M,L}=(S_L,\Quote(D_{M,L})).
\tag{10.5}\label{eq:dml}
\]

For the paper, the next theorem buys the identification of the published
\(\mathcal F\) trick, Kleene's theorem, and the IR fixed-point constructor.
\begin{theorem}[Equivalence of the three fixed-point presentations]
\label{thm:fixed-point-equivalence}
The following are translations of the same description-level construction:
\begin{enumerate}
\item the paper's explicit
\(\mathcal F(\overline{\mathcal F},\overline M,L,n,x,y,a,b)\);
\item the Kleene fixed point of the transformer which hardwires
\((\overline M,L)\) and inserts the resulting decider code into \(\Compress\);
\item the closed description \(D_{M,L}\) in \eqref{eq:dml}.
\end{enumerate}
Under the simulations of Section~\ref{sec:simulation}, their deciders compute
the same partial function.  Their description lengths are all
\(O(|M|+\log L+1)\), and their running times differ by a polynomial in
\(|M|,\log L,n\), the invoked fuel bounds, and the invoked machine times.
\end{theorem}

\begin{proof}
In the paper, Step 2 of \(\mathcal F\) applies the hardwiring operation to the
eight-input code supplied in its first argument.  When that argument is
\(\overline{\mathcal F}\), the resulting five-input code is precisely the code
of the outer decider
\(D(n,x,y,a,b)=\mathcal F(\overline{\mathcal F},\overline M,L,n,x,y,a,b)\).
Thus the string given to \(\Compress\) is its own decider description.  This
is the concrete \(s\)-\(m\)-\(n\) operation of Lemma~\ref{lem:smn} with the
diagonal argument already chosen.

Alternatively, let \(Q(e)\) output the code which performs the halting test and,
on failure, applies \(\Compress\) to \((S_L,e)\) and runs the returned decider.
Theorem~\ref{thm:kleene-two} gives \(e_Q\) with
\(\varphi_{e_Q}=\varphi_{Q(e_Q)}\), which is exactly the behavior just
described.  Equation~\eqref{eq:ycode-source} performs the same substitution
with \(e_Q\) represented by \(d_P\); Theorem~\ref{thm:ycode} proves the
operational equation for \(D_{M,L}\).  Translating it to a machine with
Theorem~\ref{thm:l-to-tm}, or translating \(\mathcal F\) with
Theorem~\ref{thm:tm-to-l}, preserves that behavior.

The displayed term contains one copy of \(M\), one binary literal \(L\), and
fixed code for the remaining operations.  Theorem~\ref{thm:ycode} therefore
gives \(|D_{M,L}|=O(|M|+\log L+1)\).  The paper's hardwiring and the size
calculation in Theorem~\ref{thm:kleene-two} give the same bound for the first
two forms.  Finally, the two simulation theorems replace each bounded call by a
fixed polynomial overhead.  Composition of finitely many such polynomials is
a polynomial in the listed quantities, proving the time statement.
\end{proof}

The paper notes that an earlier version used Kleene's recursion theorem and
that its published presentation uses \(\mathcal F\) explicitly.
\gt{fig:halt_f}{ground-truth/gt-12-compression.tex}  Theorem~\ref{thm:fixed-point-equivalence}
explains why those presentations differ in notation rather than in the
description passed to compression.

What must \(\Compress\) do with the quoted body?  It first applies
\(\Specialize\) to the static verifier and parameter holes, obtaining closed
syntax with the size bound of Lemma~\ref{lem:specialization}.  It then traverses
that syntax to build the descriptions for introspection, answer reduction, and
repetition, and it records the resulting byte lengths and fuel bounds.  It
never asks whether two terms compute the same function.  Such extensional
equality is neither available nor needed.  This self-reference is not the hard
part: it supplies no low-degree soundness, rigidity, repetition, or gap
preservation.  Those remain the substantive cited theorems.

